\subsection*{让两个查询都先结束}

先查后写很容易在顺序测试中表现正常。R1 查不到 S，购买 P1 提交；R2 随后查询，找到 P1，返回原结果。这个测试只覆盖 R1 已经完成以后 R2 才开始查找的历史。两个查询之间缺少固定顺序时，服务还有另一种执行机会。

为了明确展示它，可以让每个请求完成查找后先停下来。等到 R1 和 R2 都报告查询结束，再允许它们继续。这种安排叫并发屏障：多个执行者到达指定位置以后共同等待，等所有参与者都到齐再放行。屏障只调整请求的推进顺序，使两个空查询稳定地出现在任意写入以前；数据库仍按原有约束和提交规则处理后续语句。

从账户 1 初始余额 10000 分、空业务记录开始。R1 在连接 A 中开始事务，查询完整身份 S，得到零行。R2 在连接 B 中开始事务，也得到零行。两个查询都准确：它们各自查询时没有可见的已提交记录。两份查询结果都没有为查询者取得未来插入 S 的独占资格。

\begin{figure}[htbp]
\centering
\begin{tikzpicture}[x=1cm,y=1cm]
\node[booknote] at (2,0.35) {连接 A：R1};
\node[booknote] at (7,0.35) {共同调度};
\node[booknote] at (12,0.35) {连接 B：R2};
\node[booknode,text width=3.5cm] (a) at (2,-.45) {查找 S：零行};
\node[booknode,text width=3.5cm] (b) at (12,-.45) {查找 S：零行};
\node[bookresult,text width=3.2cm] (barrier) at (7,-1.5) {两侧到齐\\屏障放行};
\draw[bookarrow] (a.south)--++(0,-.4)-|(barrier.west);
\draw[bookarrow] (b.south)--++(0,-.4)-|(barrier.east);
\node[booknode,text width=3.5cm] (p1) at (2,-2.8) {P1、O1、D1 提交\\余额 10000 → 7000};
\node[booknode,text width=3.5cm] (p2) at (12,-2.8) {P2、O2、D2 提交\\余额 7000 → 4000};
\draw[bookarrow] (barrier.west)--(p1.north);
\draw[bookarrow] (p1.east)--node[above,font=\footnotesize]{B 随后继续扣款}(p2.west);
\node[booknote,text width=13cm] at (7,-4) {两次查询都在写入以前结束；放行以后，两条路径各自完成购买};
\end{tikzpicture}
\caption{并发屏障使两个相同身份的请求都先取得空查询结果。}
\label{fig:ch04-barrier}
\end{figure}


放开屏障后，A 生成购买编号 P1，B 生成 P2。购买编号由服务用于标识各自将要写入的购买记录，这次为两条路径选择不同编号。A 把 P1 写入购买表，把余额从 10000 扣到 7000，建立 O1 和 D1，保存成功结果并提交。B 的余额更新可能等待 A，但它仍拥有一条独立的 P2 处理路径。

B 继续以后，按数据库中的余额减去 3000，余额从 7000 变为 4000；它建立 O2 和 D2，也提交成功。最终数据库有 P1、P2 两笔购买，O1、O2 两张订单，D1、D2 两条 3000 分流水。R1 与 R2 携带同一 S，Alice 只确认一次购买，这条历史却形成了两份已提交效果。

\begin{table}[htbp]
\centering\small
\begin{tabularx}{\linewidth}{@{}>{\raggedright\arraybackslash}p{.9cm}XXp{1.8cm}@{}}
\toprule
事件 & 连接 A 处理 R1 & 连接 B 处理 R2 & 已提交余额\\
\midrule
1 & 查询 S，得到零行 & 尚未查询 & 10000\\
2 & 在屏障等待 & 查询 S，得到零行 & 10000\\
3 & 创建 P1 并扣 3000 & 等待或准备创建 P2 & 10000\\
4 & 提交 P1、O1、D1 & 继续执行 & 7000\\
5 & 已完成 & 提交 P2、O2、D2 & 4000\\
\bottomrule
\end{tabularx}
\caption{两次空查询结束以后，两条购买路径分别通过。}
\end{table}

第 3 行的已提交余额仍为 10000，因为 A 的修改尚在事务内部。A 自己可以读取扣后的 7000，其他普通查询通常仍看到先前已提交版本。表格按提交时点记录观察值，区分事务内部修改与其他连接已经能够读取的事实。

屏障使这条交错能够重复安排：两个独立连接先完成空查，再依次允许购买提交。余额更新上的等待只改变谁先扣款；没有针对 S 的约束，两条路径都仍有机会建立自己的购买。

\subsection*{原有唯一键分别保护哪个对象}

唯一约束是数据库对一列或一组列施加的规则：任何两条已经成立的记录都不能在指定列上具有相同值。购买表的 \code{purchase\_id} 唯一，能挡住第二条 P1；订单表的 \code{order\_id} 唯一，能挡住第二条 O1。给扣款流水的 \code{purchase\_id} 加唯一约束，则能限制同一购买编号至多对应一条扣款。

把这些规则用于上面的历史，P1 与 P2 不同，O1 与 O2 不同，D1 与 D2 也不同。每条流水各自只引用自己的购买编号。所有约束都可以满足，两个请求仍然重复实现了同一用户操作。已有约束观察的是服务新建的记录身份，它们还没有共同观察客户端在首次发送前保存的 S。

可以把这项差距写成两次计数。按 P1 数，成功扣款只有 D1；按 P2 数，只有 D2。按 S 数，则有 D1 和 D2 两条。要限制同一 S 的效果上限，数据库需要保存 S 与购买的关联，并对 S 本身施加唯一约束。购买表里的键值只是普通列时，数据库允许两条记录都写 K1；字段名叫 \code{idempotency\_key} 也不会改变这一点。

\begin{figure}[htbp]
\centering
\begin{tikzpicture}[x=1cm,y=1cm]
\node[bookresult,text width=3cm] (s) at (2,-1) {同一身份 S\\Alice 一次确认};
\node[booknode,text width=3.2cm] (p1) at (7.2,0) {购买 P1\\编号唯一};
\node[booknode,text width=3.2cm] (p2) at (7.2,-2) {购买 P2\\编号也唯一};
\node[booknode,text width=3.2cm] (d1) at (12,0) {O1、D1\\只引用 P1};
\node[booknode,text width=3.2cm] (d2) at (12,-2) {O2、D2\\只引用 P2};
\draw[bookalert] (s.east)--(p1.west);
\draw[bookalert] (s.east)--(p2.west);
\draw[bookarrow] (p1)--(d1);
\draw[bookarrow] (p2)--(d2);
\node[booknote,text width=13cm] at (7,-3.3) {逐笔检查：$D(P1)=D(P2)=1$\\按操作检查：$N(S)=2$，$E(S)=2$};
\end{tikzpicture}
\caption{购买编号的唯一性保护各笔记录，完整身份才把两条路径放在一起检查。}
\label{fig:ch04-existing-keys}
\end{figure}


把完整身份为 S 的已提交购买记录集合记为 \(P(S)\)，其条数记为 \(N_T(S)=|P(S)|\)；一笔购买 P 关联的已提交扣款条数记为 \(N_D(P)\)。旧规则可以逐笔满足 \(N_D(P)\leq 1\)，同时出现 \(N_T(S)=2\)。属于 S 的总扣款条数记为 \(N_D(S)\)，由各笔购买贡献的扣款相加：
\[
 N_D(S)=\sum_{P\in P(S)}N_D(P).
\]
P1 与 P2 各贡献 1，得到 \(N_D(S)=2\)。这正是第二章按业务操作统计的成功效果数；现在有了完整 S，才能从记录中确定统计范围。因此，仅把每笔购买内部的扣款上限做对，还缺少限制同一 S 关联多少笔购买的规则。事务内部执行过又回滚的 UPDATE 不计入其中。

在单个应用进程里加一把互斥锁，可以让该进程的两个处理线程轮流进入这段代码。服务有两个进程时，每个进程都能取得自己的锁。进程重新启动，内存中的锁和已处理集合也会消失。若所有写入最终落到同一个权威数据库，把竞争放进该数据库，可以使不同服务进程共同接受同一约束。

这依赖所有购买写入路径都遵守规则。后台补单、管理入口和定时恢复若绕过 S 直接新建购买，仍会产生未经仲裁的路径。权威数据库指当前对这组业务事实负责、决定哪些事务成功提交的数据库；本章只有一个这样的数据库，身份、购买和账户变化都在其中处理。多个独立数据库各自接受同一 S，已经超出这一前提。

\subsection*{查到无记录以后还缺什么}

Read Committed 的普通 SELECT 沿用一条语句开始时的快照，同时能看到本事务先前做出的修改\cite{pg}。它对当前可见记录的读取，与取得未来插入资格，是两个步骤。

有人会在查询前加 SELECT FOR UPDATE。这个语句可以锁住实际查到的行，使其他修改该行的事务等待。S 的行尚不存在时，当前例子的查询返回零行，没有一条 S 行供它取得行锁。PostgreSQL 18 的 Read Committed 也不会因为这个空查询，就普遍禁止其他事务插入满足相同条件的新行。

\begin{figure}[htbp]
\centering
\begin{tikzpicture}[x=1cm,y=1cm]
\node[booknode,text width=4.8cm] (empty) at (3,0) {按 S 查询并要求行锁\\返回零行};
\node[booknode,text width=4.8cm] (account) at (11,0) {按账户 1 查询并要求行锁\\找到已有账户行};
\node[bookresult,text width=4.8cm] (none) at (3,-1.8) {没有 S 行可锁\\其他事务仍可尝试插入 S};
\node[booknode,text width=4.8cm] (locked) at (11,-1.8) {锁住账户 1\\其他同账户修改等待};
\draw[bookarrow] (empty)--(none);
\draw[bookarrow] (account)--(locked);
\node[booknote,text width=13cm] at (7,-3.1) {PostgreSQL 18 Read Committed\\锁住实际找到的行，与约束尚未出现的 S，是不同的协调对象};
\end{tikzpicture}
\caption{空查询没有取得一条身份行的锁，预锁账户则扩大了串行范围。}
\label{fig:ch04-empty-lock}
\end{figure}


可以把这里的空结果读成一项带时点的事实：在本条语句采用的可见范围中，没有满足 S 的行。若要由它推出以后只有本请求能创建 S，还需要额外的排他规则，约束那些将在查询结束后发生的插入。单独一次读取没有写下这项规则，另一个请求便能在相同初始状态中取得同样的空结果。

有些读者可能会想到提前锁住账户 1，再查 S，使同一账户的全部购买在账户锁后串行通过。这里的账户行一直存在，因此确实有可供竞争的对象。但它会让不同 S 的身份查找也一起等待，并要求所有相关路径严格先锁账户、再做查找和写入。锁住账户本身仍未表达 S 只能对应一笔购买的数据规则。增加直接覆盖 S 的唯一约束，可以让数据库持续检查这项关系，并使遗漏协调的插入直接失败。

先查一次可以作为减少重复插入的优化：已经完成的请求命中记录后，少走一次竞争步骤。查不到的分支仍然要尝试受约束的写入。能否进入扣款步骤，取决于这次写入是否真正取得 S；先前的空查询只帮助决定去尝试。
\FloatBarrier
